\documentclass[10pt,a4paper]{amsart}
\usepackage[margin=19mm]{geometry}
\usepackage{amsmath,amssymb,mathtools}
\usepackage[scr=rsfs,cal=euler]{mathalfa}
\usepackage{xfrac}
\usepackage[unicode,hidelinks,bookmarks=false]{hyperref}
\newcommand{\ie}{i.e.\ }
\newcommand{\cf}{cf.\ }
\newcommand{\resp}{respectively\ }
\newcommand{\Subsection}{\subsection}
\providecommand{\nobreakdash}{\nobreak\hbox{-}}
\setlength{\emergencystretch}{2em}
\multlinegap0pt
\usepackage{tikz}
\usetikzlibrary{cd}
\usepackage{mathtools}
\usepackage{amssymb}
\usepackage{xcolor}
\usepackage[shortlabels]{enumitem}
\newtheorem{lemma}{Lemma}
\newcommand{\cF}{{\mathcal F}}
\newcommand{\cP}{\mathcal{P}}
\newcommand{\h}{\mathbf{h}}
\DeclareMathOperator{\op}{op}
\newcommand{\plim}[1]{\,\underset{#1}{\underset{\longleftarrow}{\lim}}\,}
\title{Interval Decompositions for Multipersistence Modules over Finite Posets and Robustness of Sheaf Data on Simplicial Complexes}
\author{Pablo Hern\'andez-Garc\'ia, Daniel Hern\'andez Serrano, Dar\'io S\'anchez G\'omez}
\date{Source: arXiv:2607.28134v1 (CC BY 4.0)}
\begin{document}
\maketitle

\noindent\emph{Source and licence.} Interval Decompositions for Multipersistence Modules over Finite Posets and Robustness of Sheaf Data on Simplicial Complexes. Version arXiv:2607.28134v1 (CC BY 4.0), distributed by arXiv under the Creative Commons Attribution 4.0 International licence, http://creativecommons.org/licenses/by/4.0/, as recorded in the arXiv licence metadata for this exact version and shown on the abstract page https://arxiv.org/abs/2607.28134v1. Full licence notice: \texttt{licenses/arxiv-2607.28134-CC-BY-4.0.txt}.\par\smallskip

\noindent\textit{Editorial scope.} Counted unit: the article's lemma Lem:DescomposicionComponentesConexas (source0.tex lines 1139--1144). The article's complete original proof of it is delivered with it (source0.tex lines 1145--1179, one \texttt{proof} environment of 35 lines). No mathematics is written, repaired or re-derived: the statement and the proof are the article's own lines, in the article's own notation, and no retained line is substituted or re-flowed.  No further source-local context is needed: every cross-reference of every retained line resolves inside the two delivered spans.  Nothing was omitted from any retained span, so the package carries no omission record.  No retained line contains a citation, so the wrapper carries no bibliography and no bibliography entry.  The retained lines contain the article's own diagram code, which the wrapper typesets with the standard tikz packages; no image file is delivered and the package declares no asset dependency.  Editorial formatting only, in addition: an AMS \texttt{amsart} preamble that restates the article's own declarations and macros used by the retained lines, namely the environments \texttt{lemma} on separate counters, together with the standard packages those lines need.  The retained environments print in this document as Lemma 1 in order of appearance, whereas the article prints the same items with the numbering of its own full document.  The complete native source of the article as distributed in the arXiv e-print for this exact version is bundled unchanged as \texttt{sources/206388/source0.tex}.
\begin{lemma}\label{Lem:DescomposicionComponentesConexas}
    Let $X_1,\dots,X_k$ be the connected components of $X$, and let $N_j\coloneqq \dim X_{j}$. Denote by $i_j$ the natural inclusion $\cP_{N_j}^{*,\op}\subseteq \cP^{*,\op}_N$. Then
    \begin{equation*}
        H^{n,\bullet}(X;\cF)\simeq {i_{1}}_*\,H^{n,\bullet}(X_1;\cF_{|_{X_1}})\,\oplus\,\dots\, \oplus\, {i_k}_*\,H^{n,\bullet}(X_k;\cF_{|_{X_k}})\, .
    \end{equation*}
\end{lemma}

\begin{proof}
    For each $\h\in \cP^{*}_N$, we have $P_X^{\h}=\bigsqcup P_{X_j}^\h$, and therefore
    \[
    H^{n,\h}(X;\cF)=H^{n}(P_X^\h;\widehat{\cF}_{|_{P_X^\h}})=\bigoplus_{j=1}^k H^n(P_{X_j}^\h;\widehat{\cF}_{|_{P_{X_j}^{\h}}})\, ,
    \]
    from which we deduce that
    \begin{equation}\label{Eq:DescomposicionComConexas}
        H^{n,\bullet}(X;\cF)\simeq \bigoplus_{j=1}^k H^{n}(P_{X_j}^\bullet;\widehat{\cF}_{|_{P_{X_j}^\bullet}})\,.
    \end{equation}
    On the other hand, for each $1\leq j\leq k$,
    \[
    \left({i_j}_* H^{n,\bullet}(X_j;\cF_{|_{X_j}})\right)_\h=
        \plim{i_j(\h')\geq_{\op} \h}H^{n,\h'}(X_j;\cF_{|_{X_j}})\, .
    \]
    The subset $\h_j\coloneqq \{h\in \h : h\leq N_j\}$ is the least element (with respect to the opposite order) such that $i_j(\h_j)\geq_{\op}\h$. Therefore, the natural morphism
    \[
    \plim{i_j(\h')\geq_{\op} \h}H^{n,\h'}(X_j;\cF_{|_{X_j}})\longrightarrow H^{n,\h_j}(X_j;\cF_{|_{X_j}})
    \]
    is an isomorphism. Moreover, for each $\h\leq_{\op} \overline{\h}$, the following diagram commutes:
    \[
    \begin{tikzcd}
        \plim{i_j(\h')\geq_{\op} \h}H^{n,\h'}(X_j;\cF_{|_{X_j}})\arrow[r,"\sim"]\arrow[d] & H^{n,\h_j}(X_j;\cF_{|_{X_j}})\arrow[d] \\
        \plim{i_j(\h')\geq_{\op} \overline{\h}}H^{n,\h'}(X_j;\cF_{|_{X_j}})\arrow[r,"\sim"] &  H^{n,\overline{\h}_j}(X_j;\cF_{|_{X_j}})
    \end{tikzcd}
    \]
    Taking into account that $P_{X_j}^{\h_j}=P_{X_j}^\h$, we deduce:
    \[
    H^{n,\h_j}(X_j;\cF_{|_{X_j}})=H^n(P_{X_j}^{\h_j};(\widehat{\cF_{|_{X_j}}})_{|_{P_{X_j}^{\h_j}}})=H^{n}(P_{X_j}^{\h_j};\widehat{\cF}_{|_{P_{X_j}^{\h_j}}})=H^n(P_{X_j}^\h;\widehat{\cF}_{|_{P_{X_j}^\h}})\,.
    \]
    That is, for each $1\leq j\leq k$, there exists an isomorphism of functors
    \[
    {i_j}_* H^{n,\bullet}(X_j;\cF_{|_{X_j}})\stackrel{\sim}\longrightarrow H^{n}(P_{X_j}^\bullet;\widehat{\cF}_{|_{P_{X_j}^\bullet}})\,, 
    \]
    Substituting into the isomorphism of Equation~\eqref{Eq:DescomposicionComConexas}, the conclusion follows.
\end{proof}

\end{document}
